\documentclass{article}
\usepackage[T1]{fontenc}
\usepackage{lmodern}
\usepackage{graphicx}
\usepackage{amsmath,amssymb}
\usepackage[a4paper,margin=2cm]{geometry}
\usepackage[absolute,overlay]{textpos}
\setlength{\TPHorizModule}{1mm}
\setlength{\TPVertModule}{1mm}
\usepackage[indonesian]{babel}
\usepackage{hyperref}
\setlength{\emergencystretch}{2em}
% unit: Halaman 44

\begin{document}

% === Halaman 44 (python/native) ===

44

1.

Inisialisasi larik S: S0 = 0, S1 = 1, …, S255 = 255 




 for i  0 to 255 do


   S[i]  i


 end 

0 

1 

2 

3 

4 

5 

6 

7 

… 




… 

252 

253 

254 

255 

0 

1 

2 

3 

4 

5 

6 

7 

… 


… 


… 

252 

253 

254 

255 

Key-Scheduling Algorithm (KSA)

\end{document}
